\documentclass[11pt,a4paper]{article}
\usepackage[utf8]{inputenc}
\usepackage[margin=0.75in]{geometry}
\usepackage{amsmath,amssymb,amsfonts}
\usepackage{xcolor}
\usepackage{listings}
\usepackage{titlesec}
\usepackage{hyperref}

% Colors
\definecolor{primaryblue}{RGB}{27, 54, 93}
\definecolor{secondaryblue}{RGB}{43, 87, 154}
\definecolor{accentgreen}{RGB}{0, 120, 0}
\definecolor{dangerred}{RGB}{180, 20, 20}
\definecolor{darkslate}{RGB}{15, 23, 42}
\definecolor{bglight}{RGB}{248, 250, 252}
\definecolor{codebg}{RGB}{241, 245, 249}
\definecolor{boxcue}{RGB}{235, 245, 255}
\definecolor{boxspeech}{RGB}{240, 253, 244}
\definecolor{boxgold}{RGB}{254, 252, 232}

\hypersetup{
    colorlinks=true,
    linkcolor=secondaryblue,
    urlcolor=secondaryblue
}

% Section styling
\titleformat{\section}{\Large\bfseries\color{primaryblue}}{\thesection}{1em}{}[\titlerule]
\titleformat{\subsection}{\large\bfseries\color{secondaryblue}}{\thesubsection}{1em}{}

% Listings configuration
\lstset{
    backgroundcolor=\color{codebg},
    basicstyle=\ttfamily\footnotesize\color{darkslate},
    keywordstyle=\color{secondaryblue}\bfseries,
    commentstyle=\color{gray}\itshape,
    stringstyle=\color{accentgreen},
    numbers=left,
    numberstyle=\tiny\color{gray},
    numbersep=6pt,
    frame=single,
    rulecolor=\color{gray!30},
    breaklines=true,
    showstringspaces=false
}

% Custom boxes using minipage & fcolorbox
\newcommand{\screenaction}[1]{%
\vspace{0.4em}\noindent
\fcolorbox{secondaryblue}{boxcue}{%
\begin{minipage}{0.965\textwidth}
\small\textbf{\color{primaryblue}[ON-SCREEN DISPLAY \& ACTION]}\\
#1
\end{minipage}}%
\vspace{0.4em}%
}

\newcommand{\speechnarration}[1]{%
\vspace{0.4em}\noindent
\fcolorbox{accentgreen!80!black}{boxspeech}{%
\begin{minipage}{0.965\textwidth}
\small\textbf{\color{accentgreen!60!black}[SPOKEN NARRATION --- Speak clearly and steadily]:}\\[0.2em]
#1
\end{minipage}}%
\vspace{0.4em}%
}

\newcommand{\insightcallout}[1]{%
\vspace{0.4em}\noindent
\fcolorbox{orange!80!black}{boxgold}{%
\begin{minipage}{0.965\textwidth}
\small\textbf{\color{orange!80!black}[THEORY VS. PRACTICALITY / COMPUTATIONAL REALITY]:}\\[0.2em]
#1
\end{minipage}}%
\vspace{0.4em}%
}

\begin{document}

% Title Header Card
\noindent
\fcolorbox{primaryblue}{primaryblue}{%
\begin{minipage}{0.965\textwidth}
\color{white}
\begin{center}
    {\LARGE \bfseries VIDEO PRESENTATION SCRIPT}\\[3pt]
    {\large Experiment 3: Sampling, Aliasing \& Whittaker-Shannon Sinc Reconstruction}\\[2pt]
    {\footnotesize Software-Based Digital Communication Laboratory $|$ Academic Year 2026}
\end{center}
\hrule \vspace{4pt}
\begin{tabular}{lll}
    \textbf{Presenter:} Apurba Maity & \textbf{Roll No:} \texttt{34900324001} & \textbf{Dept:} ECE, Cooch Behar GEC \\
    \textbf{Duration:} 4:30 -- 5:00 min & \textbf{Primary Focus:} Python Code \& Plots & \textbf{Closure:} Web Studio (30s) \\
\end{tabular}
\end{minipage}}

\vspace{0.8em}

\section*{Video Recording Checklist \& Windows Setup}
\begin{itemize}
    \item \textbf{Window 1 (Main Focus -- 75\% of Video):} VS Code / Jupyter Notebook (\texttt{Experiment\_03\_Lab\_Report.ipynb} and \texttt{experiment\_03.py}).
    \item \textbf{Window 2 (Plots Viewer):} High-resolution generated Matplotlib figures in \texttt{plots/} directory.
    \item \textbf{Window 3 (Interactive Closure -- 30s):} Browser window with \texttt{index.html} (Local / Cloudflare Simulator).
    \item \textbf{Window 4 (Closing):} GitHub Public Repository (\texttt{SecretiveCodeRunner/Digital\_Communication\_Lab}).
\end{itemize}

\vspace{0.5em}

% =========================================================================
\section{Scene-by-Scene Recording Script}
% =========================================================================

\subsection{Scene 1: Introduction \& Objectives (0:00 -- 0:40)}

\screenaction{
Show VS Code editor open to \texttt{Experiment\_03\_Lab\_Report.ipynb} on the title banner, displaying your Name, Roll Number, Department, and Experiment Title.
}

\speechnarration{
``Hello everyone and respected faculty. My name is \textbf{Apurba Maity}, Roll Number \textbf{34900324001}, from the Department of Electronics and Communication Engineering at Cooch Behar Government Engineering College.

Today, I will be presenting \textbf{Experiment 3} of our Software-Based Digital Communication Laboratory: \textbf{Sampling, Aliasing, and Whittaker-Shannon Sinc Reconstruction}.

Our core objective today is to explore the \textbf{Nyquist-Shannon Sampling Theorem} through Python numerical simulation, implement \textbf{ideal sinc interpolation} from first principles, observe \textbf{spectral aliasing} when undersampled, and most importantly, examine the critical differences between \textbf{pure theoretical mathematical assumptions and practical computational realities}.''
}

\subsection{Scene 2: Theoretical Formulation \& Signal Model (0:40 -- 1:25)}

\screenaction{
Scroll to Section 1 of the Notebook, highlighting the mathematical formulas for continuous test signal $x(t)$, Nyquist rate $f_{\text{Nyquist}}$, and the aliasing folding formula.
}

\speechnarration{
``Let us first establish our theoretical framework. We construct a deterministic continuous two-tone signal:
\begin{equation*}
    x(t) = \sin(2\pi \cdot 100 t) + 0.5 \sin(2\pi \cdot 300 t)
\end{equation*}
Here, the fundamental frequency is $f_1 = 100\text{ Hz}$ and the highest harmonic is $f_2 = 300\text{ Hz}$. Thus, the maximum signal bandwidth is $f_{\max} = 300\text{ Hz}$.

According to the \textbf{Nyquist-Shannon Sampling Theorem}, exact lossless recovery requires a sampling frequency satisfying:
\begin{equation*}
    f_s \ge 2 f_{\max} = 2 \times 300\text{ Hz} = 600\text{ Hz}
\end{equation*}
Continuous-time reconstruction is governed by the \textbf{Whittaker-Shannon sinc interpolation formula}:
\begin{equation*}
    \hat{x}(t) = \sum_{n=-\infty}^{\infty} x[n] \cdot \operatorname{sinc}\left(f_s (t - n T_s)\right), \quad \text{where } \operatorname{sinc}(u) = \frac{\sin(\pi u)}{\pi u}
\end{equation*}
If we violate this by undersampling at $f_s = 400\text{ Hz}$, the high-frequency $300\text{ Hz}$ component folds across the Nyquist boundary $f_s/2 = 200\text{ Hz}$ according to the aliasing formula:
\begin{equation*}
    f_{\text{alias}} = |f_2 - k \cdot f_s| = |300 - 1 \times 400| = |-100| = 100\text{ Hz}
\end{equation*}
Now, let us examine how we implemented this in Python without using external black-box toolboxes.''
}

\newpage

\subsection{Scene 3: Python Code Architecture \& DSP Walkthrough (1:25 -- 2:15)}

\screenaction{
Switch to \texttt{experiment\_03.py} in VS Code. Highlight the \texttt{sinc\_reconstruction()} function (lines 53--63) and \texttt{calculate\_spectrum()} (lines 65--73).
}

\begin{lstlisting}[language=Python,caption={Core Vectorized Sinc Reconstruction Algorithm in Python}]
def sinc_reconstruction(sample_times, sample_values, t_fine, fs):
    # Vectorized 2D broadcast: shape (N_samples, M_eval_points)
    dt_matrix = t_fine[None, :] - sample_times[:, None]
    sinc_matrix = np.sinc(fs * dt_matrix)  # np.sinc computes sin(pi*x)/(pi*x)
    # Sum overlapping sinc pulses across all discrete samples
    x_hat = np.dot(sample_values, sinc_matrix)
    return x_hat
\end{lstlisting}

\speechnarration{
``Here in our Python implementation, we avoided black-box DSP functions and implemented the core mathematics directly using NumPy.

As shown in lines 53 to 63 in \texttt{sinc\_reconstruction}:
\begin{itemize}
    \item To evaluate the Whittaker-Shannon summation efficiently for thousands of time points, we construct a 2D distance matrix using array broadcasting: \texttt{dt\_matrix = t\_fine[None, :] - sample\_times[:, None]}.
    \item Passing this into \texttt{np.sinc(fs * dt\_matrix)} creates our full interpolation kernel grid.
    \item A single vectorized dot product \texttt{np.dot(sample\_values, sinc\_matrix)} instantly sums all overlapping sinc pulses across the entire continuous evaluation grid in a fraction of a millisecond.
\end{itemize}

For frequency analysis, lines 65 to 73 use \texttt{np.fft.rfft} to compute the single-sided discrete Fourier spectrum, normalized by $2/N$ to obtain exact physical peak amplitudes.

Furthermore, in lines 110 to 125, we implemented sample boundary padding ($N_{\text{pad}} = 20$), eliminating artificial edge-truncation errors in our numerical reconstruction.''
}

\subsection{Scene 4: Plot Walkthrough \& Empirical Validation (2:15 -- 3:25)}

\screenaction{
Display the 4 generated Matplotlib figures from \texttt{plots/} or scroll through Section 3 of the Jupyter Notebook.
}

\speechnarration{
``Let us now examine our empirical results across the 4 generated figures:

\textbf{1. Discrete Sample Stems (Figure 1):} Over our 50 ms test duration:
\begin{itemize}
    \item At \textbf{1800 Hz} (top), 91 dense samples capture 6 samples per cycle of the 300 Hz wave.
    \item At \textbf{600 Hz} (middle), 31 samples capture exactly 2 samples per cycle.
    \item At \textbf{400 Hz} (bottom), only 21 samples are recorded---barely 1.3 samples per cycle---missing the sinusoidal peaks.
\end{itemize}

\textbf{2. Reconstructed Waveforms \& FFT Spectra (Figures 2 \& 3):}
\begin{itemize}
    \item In the \textbf{1800 Hz Oversampled case}, the reconstructed dashed curve overlays the original black reference signal perfectly, yielding a near-zero Mean Squared Error of $5.0 \times 10^{-6}\text{ V}^2$. The FFT spectrum cleanly recovers both the 100 Hz and 300 Hz peaks.
    \item In the \textbf{400 Hz Undersampled case}, the reconstructed waveform is heavily distorted. Looking at the FFT spectrum, the 300 Hz peak has completely vanished from its true position and folded directly onto \textbf{100 Hz}, exactly confirming our theoretical derivation $|300 - 400| = 100\text{ Hz}$!
\end{itemize}

\textbf{3. Time-Domain Error (Figure 4):} The error waveform $e(t) = x(t) - \hat{x}(t)$ confirms this quantitatively: oversampling yields a flat zero line, while undersampling results in a massive oscillating error with $\text{MSE} = \mathbf{0.25\text{ V}^2}$.''
}

\newpage

\subsection{Scene 5: Theory vs. Practicality \& Computational Realities (3:25 -- 4:15)}

\screenaction{
Scroll to \textbf{Section 4: Theory vs. Practicality} in the Notebook. Point with mouse cursor to the comparative table and error numbers.
}

\insightcallout{
\textbf{Key Insight:} Explain why $f_s = 600\text{ Hz}$ produced $\text{MSE} = 0.1235$ despite being at the Nyquist rate, and discuss finite-window sinc truncation and hardware anti-aliasing filters.
}

\speechnarration{
``Now, let us address the most insightful question of this experiment: \textbf{Where does pure mathematical theory diverge from practical simulation and physical hardware?}

\textbf{1. Phase Nulling at Critical Nyquist ($f_s = 2f_{\max} = 600\text{ Hz}$):}
In textbooks, theory states that $f_s = 2f_{\max}$ guarantees exact signal reconstruction. However, in our simulation, the 600 Hz critical case produced an MSE of \textbf{0.1235}! Why did this happen?

Our test signal is a pure sine wave: $\sin(2\pi \cdot 300 t)$. When sampled at intervals of $T_s = 1/600\text{ s}$, every sample falls exactly on a zero crossing:
\begin{equation*}
    x_{\text{harmonic}}[n] = 0.5 \sin\left(2\pi \cdot 300 \cdot \frac{n}{600}\right) = 0.5 \sin(n\pi) \equiv 0
\end{equation*}
Every single sample of the 300 Hz tone evaluated to zero! The sinc interpolator was therefore blind to the 300 Hz wave and only reconstructed the 100 Hz tone. The missing harmonic power is exactly $\frac{A_2^2}{2} = \frac{0.5^2}{2} = \mathbf{0.125\text{ V}^2}$, which matches our measured MSE of $0.1235\text{ V}^2$!

This proves that in real-world DSP and ADC design, we cannot operate right at the theoretical boundary; we must strictly oversample ($f_s \ge 2.2 f_{\max}$) with a safety guard band to avoid phase nulling.

\textbf{2. Infinite Sinc Summation vs. Finite Observation Window:}
Theoretical Whittaker-Shannon interpolation requires an infinite summation from $-\infty$ to $+\infty$. Real DSP processors and digital scopes operate on finite time windows ($T = 50\text{ ms}$). Truncating the sinc series introduces boundary ripples, which we resolved by implementing sample padding ($N_{\text{pad}} = 20$).

\textbf{3. Physical Anti-Aliasing Filtering (AAF):}
In software math, we can inspect and predict folded frequencies. But in physical ADC hardware, once an out-of-band high-frequency signal aliases into baseband, it is mathematically inseparable from legitimate data. Therefore, an analog active Low-Pass Anti-Aliasing Filter must always precede the hardware ADC.''
}

\subsection{Scene 6: Interactive Web Simulator (Bonus Closure) (4:15 -- 4:45)}

\screenaction{
Switch to browser window displaying \texttt{index.html}. Click preset buttons (\textbf{1800 Hz}, \textbf{600 Hz}, \textbf{400 Hz}) and briefly drag the $f_s$ slider.
}

\speechnarration{
``As a bonus interactive closure, we also compiled this lightweight 60 FPS HTML5 simulation studio.

Here, we can dynamically drag the sampling frequency slider from 1800 Hz down to 300 Hz. Notice how the reconstructed green waveform morphs in real time, and in the live FFT spectrum below, you can watch the aliased red peak sweep continuously across the folding boundary as $f_s$ decreases. We can also synthesize Web Audio tones to hear how aliasing lowers the perceived auditory pitch.''
}

\subsection{Scene 7: Summary \& Conclusion (4:45 -- 5:00)}

\screenaction{
Switch to GitHub repository page (\texttt{SecretiveCodeRunner/Digital\_Communication\_Lab}).
}

\speechnarration{
``In conclusion, this experiment successfully verified the Nyquist-Shannon Sampling Theorem, implemented vectorized sinc reconstruction in Python, validated spectral aliasing folding, and clarified the boundary between continuous theory and discrete DSP implementation.

The complete Python code, Jupyter lab notebook, printable study guides, and the interactive studio are published on my open-source GitHub repository.

Thank you very much for your time and attention!''
}

\newpage

% =========================================================================
\section{Summary Data Table \& Viva Q\&A Defense Sheet}
% =========================================================================

\subsection{Empirical Simulation Summary Table}

\begin{center}
\renewcommand{\arraystretch}{1.3}
\begin{tabular}{|l|c|c|c|c|c|l|}
\hline
\textbf{Regime} & \textbf{$f_s$ (Hz)} & \textbf{Ratio} & \textbf{Samples} & \textbf{MSE ($\text{V}^2$)} & \textbf{RMSE ($\text{V}$)} & \textbf{Harmonic Status} \\
\hline
\textbf{Oversampled} & 1800 & $3.0\times$ & 91 & $5.0 \times 10^{-6}$ & 0.0022 & Clean recovery at 300 Hz \\
\hline
\textbf{Critical} & 600 & $1.0\times$ & 31 & 0.1235 & 0.3514 & Phase Null ($\sin(n\pi)=0$) \\
\hline
\textbf{Undersampled} & 400 & $0.67\times$ & 21 & 0.2500 & 0.5000 & \textbf{Aliased to 100 Hz} \\
\hline
\end{tabular}
\end{center}

\vspace{1em}

\subsection{Faculty Viva \& Defense Cheat-Sheet}

\noindent
\fcolorbox{secondaryblue}{bglight}{%
\begin{minipage}{0.965\textwidth}
\small
\textbf{\color{primaryblue}Q1: Why did Critical Sampling ($f_s = 600\text{ Hz}$) yield an MSE of 0.1235 instead of 0?}\\[0.3em]
\textbf{Answer:} The 300 Hz harmonic was defined as $0.5\sin(2\pi \cdot 300 t)$. When sampled at intervals of $T_s = 1/600\text{ s}$, the sample times are $t_n = n/600$, resulting in $0.5\sin(n\pi) \equiv 0$ for all integers $n$. Thus, all discrete samples of the 300 Hz component were identically zero! Sinc interpolation could only recover the 100 Hz tone. The energy of the lost 300 Hz tone is $\frac{1}{2} A_2^2 = \frac{1}{2}(0.5)^2 = \mathbf{0.125\text{ V}^2}$, which matches our measured MSE of $0.1235\text{ V}^2$.
\end{minipage}}

\vspace{0.6em}
\noindent
\fcolorbox{secondaryblue}{bglight}{%
\begin{minipage}{0.965\textwidth}
\small
\textbf{\color{primaryblue}Q2: How does your Python code implement sinc reconstruction without slow loops?}\\[0.3em]
\textbf{Answer:} Through 2D array broadcasting and matrix multiplication. We construct an $(M \times N)$ distance matrix \texttt{dt = t\_fine[None, :] - sample\_times[:, None]}, apply \texttt{np.sinc(fs * dt)}, and execute \texttt{np.dot(sample\_values, sinc\_matrix)}. This vectorizes thousands of overlapping sinc evaluations in a single C-optimized NumPy call.
\end{minipage}}

\vspace{0.6em}
\noindent
\fcolorbox{secondaryblue}{bglight}{%
\begin{minipage}{0.965\textwidth}
\small
\textbf{\color{primaryblue}Q3: Why is an Anti-Aliasing Filter essential in hardware ADCs?}\\[0.3em]
\textbf{Answer:} Once an out-of-band analog frequency higher than $f_s/2$ enters the ADC sampler, it aliases into the baseband $[0, f_s/2]$ and produces identical discrete samples as a legitimate low-frequency signal. It is mathematically impossible to filter out or distinguish aliased components after digitization. Therefore, an analog active Low-Pass Filter must attenuate out-of-band frequencies \textit{before} the ADC.
\end{minipage}}

\vspace{0.6em}
\noindent
\fcolorbox{secondaryblue}{bglight}{%
\begin{minipage}{0.965\textwidth}
\small
\textbf{\color{primaryblue}Q4: Why did we pad samples ($N_{\text{pad}} = 20$) in the time domain?}\\[0.3em]
\textbf{Answer:} The theoretical Whittaker-Shannon formula requires an infinite summation ($-\infty \le n \le +\infty$). In numerical simulation over a finite window $[0, T]$, truncating the sinc series leaves points near $t = 0$ and $t = T$ without neighboring sinc pulses, causing boundary edge distortion. Evaluating samples slightly beyond the plotting window ($[-20 T_s, T + 20 T_s]$) ensures full sinc overlap across the entire evaluation interval.
\end{minipage}}

\vfill
\begin{center}
    \footnotesize \color{gray}
    Document compiled with \LaTeX\ for Apurba Maity $|$ Cooch Behar Government Engineering College
\end{center}

\end{document}
